\section{问题重述与分析}
\subsection{问题背景}
非孤岛式微网由外部电网、光伏发电和储能设备共同向小区供电。光伏集中于白天，负载和电价则具有各自的日内变化规律。储能能够转移供电时段，但充放电损耗、功率上限和储电量边界限制了这种转移。因此，购电策略既要满足当前负载，也要为后续时段保留适当的储能，不能只依据当时的供需差额逐段决定。

\subsection{问题要求}
题目给定储能额定容量12\,000 kWh、最大充放电功率5\,000 kW及充放电效率90\%，要求储电量保持在1\,200--10\,800 kWh之间。结合单日预测、全年实际数据和定时发布的光伏预报，需要依次解决以下问题。
\begin{enumerate}
  \item 在已知当日电价、负载和光伏预测的条件下，于0:00确定全天计划购电量与充放电量，保证供电并满足日初、日末储电量相等。
  \item 当实际负载与光伏变化时，根据日初信息制定购电计划；供电不足部分按交易时刻电价的5倍紧急购入，计划费用仍按原计划购电量计算。
  \item 利用0、6、12、18时发布的光伏预报制定和调整计划，并计入减购违约费、增购费及紧急购电费，讨论其他预报时刻的必要性。
  \item 将上述后两类合同条件推广到波动电价，重新计算相应购电与储能方案，输出指定日期结果及完整策略文件。
\end{enumerate}

\subsection{建模思路}
问题一中，全天输入在日初给定，主要困难是储能状态对各时段的连接。将功率换算为区间电量后，供需平衡和储能递推均为线性关系；再用二进制变量表示充放电互斥，即可在统一时域内联合求解144个区间的决策。

问题二中，日内正常购电合同在0:00冻结，未来实际供需只能由可见历史预测。故先在24小时滚动窗口内安排合同和储能动作，再以实际供需回放确定紧急购电与费用；合同计划、预测目标和实际账单分别保存，避免将三者混用。

问题四中，未来负载、光伏和电价不能作为已知量直接使用。先根据决策时刻已获得的信息形成点预测，再以已实现的联合误差构造采购情景，比较正常合同购电与高价缺口补电的费用，在合同权限内优化共同购电和储能动作。实际供需偏差由本地动作保护及紧急购电处理，实际后继储电量成为下一次优化的初态。费用按实际合同版本和交易价格逐项结算，情景规划目标与实际支出分别核对。

本文完成问题一、问题二和问题四的建模、结果分析及检验；问题三保留待补充章节。
